\section{Introduction}
\label{sec:introduction}

Python native extensions connect Python programs to C and C++ code. This
mechanism supports numerical kernels, device runtimes, operating-system
interfaces, and other functionality that cannot be implemented efficiently in
pure Python. The same boundary removes native code from the memory and error
management provided by the interpreter. Extension developers must instead
follow the Python/C API rules for reference ownership, nullable returns,
exception state, allocator domains, buffers, garbage collection, and thread
state \cite{pythoncapi,hu2020pythoncapi}. A violation can leak an object, access
a released value, corrupt interpreter state, or terminate the process that
hosts the extension.

These rules are difficult to enforce for two reasons. First, the relevant
information is distributed across the Python/C API reference manual. The
correct action depends on details such as whether an API returns a new or
borrowed reference, steals an argument, preserves an exception, or pairs with
a specific release operation. Second, compliance is path-sensitive. A newly
created object can be returned on one path, transferred to a container on a
second path, and require an explicit decrement on every remaining path. A
cleanup label that is correct for an owned reference can release a borrowed
reference on another branch. The defect is therefore not characterized by an
isolated call or a fixed API pair.

Prior Python/C analyzers make important progress on reference-count errors.
Pungi transforms reference-count behavior into affine programs
\cite{pungi2014}, while PyRefcon tracks both reference counts and object
lifecycle states with symbolic execution \cite{pyrefcon2023}. These analyses
target reference management and require detailed models of API and program
behavior. They do not provide a common representation for the broader set of
Python/C API requirements, including error-state discipline, allocator
agreement, buffer lifetime, numeric conversion, garbage-collection protocols,
and capsule metadata. General native-code analyzers cover some low-level
defects but lack the ownership and interpreter-state semantics that determine
whether a Python/C API use is correct.

Large language models offer a complementary ability to reason about source
code and natural-language specifications. Unconstrained bug finding, however,
asks a model to infer the relevant property, identify the program object, trace
the paths, and judge the violation in one step. This formulation makes reports
difficult to audit and exposes the result to unsupported bug hypotheses.
Recent hybrid systems improve practical static analysis with model reasoning
\cite{li2024llmstatic}, and specification-guided systems decompose external
requirements before checking code \cite{safe4u2025,rfcaudit2025}. The open
question for Python native extensions is how to formulate the API-specific
checking task so that static analysis and model reasoning have separate,
auditable roles. Unlike a typestate checker, \tool does not build a complete
transfer function for the whole program; unlike a pattern-based checker,
obligation definitions are separated from syntactic suspicion; and unlike
unconstrained LLM bug finding, the model is constrained to a fixed
API-induced property with enumerated safe alternatives.

We make the observation that a Python/C API requirement can be represented as
a \emph{safety obligation} together with explicit \emph{discharge
conditions}. An API use creates an obligation for a program object or runtime
state. The obligation is safe only when every relevant path satisfies an
admissible discharge condition. For example, a new reference satisfies its
ownership obligation when the path releases it, transfers it to an API that
steals ownership, or returns it to the caller. This representation states what
must hold without prescribing one program-analysis technique.

Based on this observation, we develop \tool. An offline construction process
extracts normative requirements and correct usage alternatives from CPython
documentation and real bug cases. The resulting specification contains 13
obligation categories and 41 discharge conditions. At detection time,
lightweight static analysis finds Python/C API call sites, binds them to
obligations, and collects the associated program objects. Fifteen syntactic
patterns implemented by the static analyzer prune obligation instances that
lack a suspicious shape. These patterns decide which obligation instances reach
the LLM, trading recall risk for cost reduction; they are not part of the
obligation semantics. For each retained instance, a large language
model receives the enclosing function, the obligation, and its discharge conditions.
The model identifies the trigger, enumerates relevant exits, checks each
condition, and returns a structured verdict and evidence.

We evaluate \tool through three research questions. The in-the-wild study
audits 104 high-confidence reports from widely used Python extension projects.
Eighty-six reports are real previously unknown defects, yielding 82.7\%
precision, and 45 findings were confirmed by project maintainers. The baseline
study reproduces PyRefcon and audits 80 deduplicated findings; it analyzes
why PyRefcon achieves only 7.5\% precision and examines the five
PyRefcon-only true positives to expose the boundary between Python/C API
obligations and generic C memory safety. Finally, the design study measures
the cost reduction from static pattern-based pruning and ablates the
chain-of-thought discharge step and the discharge-condition list to quantify
their individual contributions to bug recall.

This paper makes the following contributions:

\begin{itemize}
  \item We formulate Python/C API correctness as dischargeable safety
  obligations and construct a specification with 13 obligation categories
  and 41 discharge conditions.
  \item We design and implement \tool, which combines static obligation
  binding and pattern-based pruning with structured, path-oriented
  LLM checking and evidence-bearing reports.
  \item We evaluate \tool on real Python native extensions, find 86 previously
  unknown bugs at 82.7\% precision, compare with the closest reproducible
  analyzer, and isolate the effects of pruning and structured checking.
\end{itemize}
